\chapter{Organizaciones de mantenimiento: Parte 145}\label{ch:18}
\objetivos{Relacionar alcance de aprobación, recursos, ejecución y certificación; diferenciar cumplimiento y gestión de seguridad.}
\section{Una aprobación tiene alcance}
Parte 145 establece requisitos para la aprobación de organizaciones de mantenimiento dentro de su ámbito. El certificado, las habilitaciones, limitaciones, ubicaciones y procedimientos definen lo que la organización puede hacer. Tener una aprobación no autoriza cualquier tarea en cualquier producto o lugar.

La autoridad competente se determina en 145.1. El punto 145.A.10 establece el ámbito; los recursos y procedimientos se desarrollan en los apartados siguientes. El alcance de esta explicación es académico y no sustituye el manual de organización ni su aprobación.
\section{Instalaciones, herramientas y componentes}
Las instalaciones deben ser adecuadas al trabajo: acceso, protección frente al ambiente, zonas y almacenamiento influyen en la ejecución. Las herramientas y equipos necesitan disponibilidad y control conforme al procedimiento, incluida calibración cuando corresponda. Una llave que muestra una cifra no ofrece por sí sola trazabilidad metrológica.

Los componentes se identifican y se gestionan según su estado. Mezclar material servible, no servible y retirado de servicio puede provocar instalaciones incorrectas. El embalaje, la preservación, la caducidad y los documentos de origen forman parte del control, no solo el número de parte.
\section{Personal y datos de mantenimiento}
145.A.30 y 145.A.35 tratan requisitos de personal y personal certificador/de apoyo. La competencia se relaciona con función y tareas. Una licencia y una autorización de certificación de la organización no son documentos idénticos; los privilegios se ejercen bajo condiciones específicas.

145.A.45 exige el uso de datos de mantenimiento aplicables y actualizados en el contexto que regula. Antes de ejecutar una tarea se identifica configuración, revisión y procedimiento. Un documento encontrado en internet con un título parecido no garantiza aplicabilidad. Si una instrucción es ambigua o incompleta se activa el procedimiento de aclaración y comunicación, no se improvisa una reparación.
\section{Planificación y ejecución}
La planificación considera personal, herramientas, repuestos, datos, instalaciones y condiciones de ejecución. Las tareas críticas pueden requerir métodos de captura de errores y verificaciones específicas. El control de herramientas y objetos extraños evita introducir daño durante una intervención.

El relevo debe comunicar trabajo completado, estado de la aeronave, elementos abiertos y acciones pendientes. Un sistema que depende solo de la memoria de una persona es vulnerable a interrupciones. La presión de producción no elimina las condiciones de seguridad y certificación.
\ilustra{parte145}{Relación funcional entre alcance aprobado, planificación, recursos, ejecución y certificación. La gestión y el seguimiento de cumplimiento actúan sobre el conjunto.}
\section{Certificación de mantenimiento y registros}
145.A.50 establece condiciones para el certificado de puesta en servicio: personal debidamente autorizado verifica que el mantenimiento ordenado se ha realizado correctamente conforme a los procedimientos y datos pertinentes y que no hay incumplimientos conocidos que comprometan seguridad de vuelo. El texto contiene supuestos y condiciones específicos; no se reduce a una casilla genérica de «todo revisado».

El trabajo en componentes fuera de la aeronave puede dar lugar a EASA Form 1 según el apartado y excepciones aplicables. Este documento no demuestra por sí solo que el componente sea instalable en cualquier aeronave: se comprueban identidad, configuración, condiciones y elegibilidad de instalación. Los registros permiten reconstruir qué se certificó y deben conservarse conforme a las reglas aplicables.
\section{Gestión de seguridad y cumplimiento}
145.A.200 exige un sistema de gestión con responsabilidades, política de seguridad, identificación de peligros, evaluación y gestión de riesgos, competencia, procesos documentados y seguimiento del cumplimiento. La edición actual no debe describirse solo mediante el antiguo esquema de «sistema de calidad» sin atender a su evolución.

El seguimiento del cumplimiento comprueba requisitos y procedimientos; la gestión de seguridad busca peligros y eficacia de mitigaciones. Pueden estar integrados, pero una auditoría de formularios no agota el análisis de riesgo. El punto 145.A.200A introduce requisitos de gestión de seguridad de la información en su ámbito y con las disposiciones aplicables; la información técnica y los sistemas digitales también pueden afectar seguridad.
\caso{M18. Trabajo completado, dato incorrecto}{Una orden ficticia está firmada, pero se detecta que se aplicó una revisión de manual no válida para la configuración. No basta con corregir la fecha del registro: se evalúa la diferencia técnica, el trabajo realmente ejecutado y las acciones necesarias por el proceso de la organización. También se investiga por qué el sistema permitió seleccionar el dato incorrecto. La respuesta combina producto, documentación y causa organizativa.}
\section{Una organización que aprende}
Los hallazgos deben llevar a acciones correctoras y a comprobar su eficacia. Cambiar un formulario no demuestra que se haya eliminado la causa. Indicadores de repeticiones, errores y uso de datos ayudan a evaluar el sistema, siempre con definiciones y exposición conocidas.

Para el estudiante, el objetivo es poder seguir la cadena desde aprobación y datos hasta tarea y certificado. Esa comprensión evita confundir capacidades técnicas de una persona con privilegios de una organización o considerar los registros como un trámite posterior sin relación con la seguridad.

\section{Relevo y captura de errores}
Un relevo identifica pasos completados, elementos abiertos, pruebas pendientes, configuración temporal y anomalías. Permite reconstruir el estado sin depender de memoria o conversación. Interrupciones y turnos son condiciones previsibles que los procedimientos y registros deben soportar.

Una segunda firma sin independencia ni evidencia de qué se comprobó puede ser solo apariencia de barrera. El método de captura debe responder al error y a los requisitos de la tarea. La competencia incluye entender el objetivo de la comprobación.
\section{Herramientas y cadena de medida}
Se necesita equipo adecuado al rango y tolerancia, estado de calibración y condiciones de uso. Calibración no declara correcta cualquier lectura: incertidumbre, resolución y técnica también influyen. Si aparece equipo fuera de tolerancia, se evalúan trabajos afectados; cambiar su etiqueta no resuelve medidas anteriores.

El control de herramientas y objetos extraños relaciona inventario, acceso, uso y cierre. Una pérdida puede requerir investigación por el posible efecto sobre producto. Almacenamiento y preservación conservan condiciones de piezas y consumibles, incluida caducidad cuando proceda.
\section{Contener el problema y eliminar su causa}
La acción inmediata contiene; la correctora busca evitar repetición. Retirar una copia incorrecta de manual puede ser necesario e insuficiente: se revisa distribución, selección, alertas y comprobación. La eficacia se demuestra después, no por la fecha de cierre del formulario.

145.A.200 relaciona seguridad y seguimiento de cumplimiento con funciones específicas. Investigar solo quién cometió el último error no explica por qué las barreras permitieron el evento. Esta perspectiva no impone una metodología única; enseña a estudiar condiciones y sistema.
\repaso{¿Un útil calibrado demuestra por sí solo apriete correcto?}{También importan selección, configuración, técnica, dato, secuencia y condiciones. Calibración es una evidencia dentro de una cadena más amplia.}

\resumen{La aprobación tiene límites. Personal, herramientas y datos deben ser adecuados a la tarea. CRS y Form 1 tienen funciones y condiciones distintas. La certificación no corrige una ejecución con datos inadecuados. Seguridad, cumplimiento e información requieren una gestión coherente.}
\fuentes{Reglamento 1321/2014, 145.1 y 145.A.10, .25, .30, .35, .40, .42, .45, .47, .48, .50, .55, .70 y .200--.200A \cite{EU1321}; documentación docente de organizaciones.}
